\documentclass[10pt,a4paper]{amsart}
\usepackage[margin=19mm]{geometry}
\usepackage{amsmath,amssymb,mathtools}
\usepackage[scr=rsfs,cal=euler]{mathalfa}
\usepackage{xfrac}
\usepackage[unicode,hidelinks,bookmarks=false]{hyperref}
\newcommand{\ie}{i.e.\ }
\newcommand{\cf}{cf.\ }
\newcommand{\resp}{respectively\ }
\newcommand{\Subsection}{\subsection}
\providecommand{\nobreakdash}{\nobreak\hbox{-}}
\setlength{\emergencystretch}{2em}
\multlinegap0pt
\usepackage{mdframed}
\usepackage{mdframed}
\usepackage{mdframed}
\usepackage{mdframed}
\usepackage{xcolor}
\usepackage{amssymb}
\usepackage{csquotes}
\usepackage{float}
\usepackage[shortlabels]{enumitem}
\usepackage{mdframed}
\usepackage{dsfont}
\usepackage{bm}
\newtheorem{lemma}{Lemma}
\title{Mean-field imitation dynamics on fast assortative networks}
\author{Benedict Russell, Andrew Nugent, Jacques Bara}
\date{Source: arXiv:2606.11892v1 (CC BY 4.0)}
\begin{document}
\maketitle

\noindent\emph{Source and licence.} Mean-field imitation dynamics on fast assortative networks. Version arXiv:2606.11892v1 (CC BY 4.0), distributed by arXiv under the Creative Commons Attribution 4.0 International licence, http://creativecommons.org/licenses/by/4.0/, as recorded in the arXiv licence metadata for this exact version and shown on the abstract page https://arxiv.org/abs/2606.11892v1. Full licence notice: \texttt{licenses/arxiv-2606.11892-CC-BY-4.0.txt}.\par\smallskip

\noindent\textit{Editorial scope.} Counted unit: the article's lemma below (source0.tex lines 243--248). The article's complete original proof of it is delivered with it (source0.tex lines 917--947, one \texttt{proof} environment of 31 lines, which the article places away from the statement). No mathematics is written, repaired or re-derived: the statement and the proof are the article's own lines, in the article's own notation, and no retained line is substituted or re-flowed.  No further source-local context is needed: every cross-reference of every retained line resolves inside the two delivered spans.  Nothing was omitted from any retained span, so the package carries no omission record.  No retained line contains a citation, so the wrapper carries no bibliography and no bibliography entry.  No retained line contains a figure environment, an image-inclusion command or drawing code, so no image is delivered and the package declares no asset dependency.  Editorial formatting only, in addition: an AMS \texttt{amsart} preamble that restates the article's own declarations and macros used by the retained lines, namely the environments \texttt{lemma} on separate counters, together with the standard packages those lines need.  The retained environments print in this document as Lemma 1 in order of appearance, whereas the article prints the same items with the numbering of its own full document.  The complete native source of the article as distributed in the arXiv e-print for this exact version is bundled unchanged as \texttt{sources/202580/source0.tex}.
\begin{lemma}\label{lem: mean bounded below}
    On any finite interval $[0,T]$, the mean is bounded below. Specifically,
    \begin{align*}
        m_\mu(t) \geq m_{\mu_0}e^{-{t}}.
    \end{align*}
\end{lemma}

\begin{proof}
    The density along characteristics satisfies
    \begin{align*}
        \frac{d}{dt}\mu(t,X(t,x)) &= - \mu(t,X(t,x)) \partial_x V_\mu(t,X(t,x)),
    \end{align*}
    and therefore
    \begin{align*}
        \mu(t,X(t,x)) = \mu_0\; \exp\Big(-\int_0^t\partial_x V_\mu(s,X(s,x))ds \bigg).
    \end{align*}
    Taking the maximum value, 
    \begin{align*}
        \|\mu(t,\cdot)\|_\infty \leq \|\mu_0\|_\infty \; \exp\Big(\int_0^t \|\partial_x V_\mu(s,\cdot)\|_\infty ds \bigg).
    \end{align*}
    By Lemma \ref{lem: mean bounded below}, the mean is bounded below on the finite time interval. Now consider an upper bound on $\partial_xV_\mu(t,x)$. Note that $I_\mu$ and $V_\mu$ are both differentiable in $x$ when $\mu \in L^\infty(\Omega)$.
    \begin{align*}
        \sup_x |\partial_x V_\mu(t,x)| &= \sup_x|\partial_x \frac{1}{x+m_\mu}I_\mu(t,x)|\\
        &= \sup_x|-\frac{I(x,\mu)}{(x+m_\mu)^2} + \frac{1}{x+m_\mu}\partial_xI(x,\mu)|\\
        &\leq |\frac{1}{m_\mu^2} + \frac{1}{m_\mu}\partial_xI_\mu(t,x)|\\
        &\leq \frac{1}{m_\mu^2} + \frac{1}{m_\mu}|\partial_x\int(y^2-x^2)\mu(t,y)p_\mu(x,y)dy|\\
        & = \frac{1}{m_\mu^2} + \frac{1}{m_\mu}|\int(-2x)\mu(t,y)p_\mu(x,y) + (y^2-x^2)\mu(t,y)\partial_xp_\mu(x,y)\; dy|\\
        &\leq \frac{1}{m_\mu^2} + \frac{1}{m_\mu}|2 + \int_\Omega (y^2-x^2)\mu(t,y)\Big(-\beta \int(bz-c)\mu(t,z) dz\; p_\mu(x,y)\\&\qquad \times (1-p_\mu(x,y)\Big)\; dy|\\
        &\leq \frac{1}{m_\mu^2} + \frac{1}{m_\mu}|2 -\frac{\beta}{4} \int(bz-c)\mu(t,z)dz|\\
        &\leq \frac{1}{m_\mu^2} + \frac{2 + \frac{\beta}{4} |\int(bz-c)\mu(t,z)dz|}{m_\mu}\\
        &\leq \frac{1}{m_\mu^2} + \frac{2 + \frac{\beta}{4} b m_\mu + c}{m_\mu}\\
        &\leq \frac{1}{\delta_T^2} + \frac{2+c}{\delta_T} + \frac{\beta b}{4}.
    \end{align*}
    The result then follows substituting the upper bound into the growth bound
    \begin{align*}
        \|\mu(t,\cdot)\|_\infty \leq \|\mu_0\|_\infty \; \exp\bigg[t\big(\frac{1}{\delta_T^2} + \frac{2+c}{\delta_T} + \frac{\beta b}{4}\big)\bigg].
    \end{align*}
\end{proof}

\end{document}
